\documentclass[11pt, letterpaper]{ctexart}
% Phase 12A — 暑期科研实习终期报告（中文版）
% 与 report/main.tex 英文版结构一致：单栏、Times New Roman + Songti SC、
% 1.5 倍行距、阿拉伯数字章节编号、封面页、目录、数字编号引用。

\usepackage{amsmath,amssymb,amsfonts}
\usepackage{graphicx}
\usepackage{textcomp}
\usepackage{xcolor}
\usepackage{booktabs}
\usepackage{tabularx}
\usepackage{caption}
% 表格题注与表体之间留更多间距；表格题注左右各缩进 3cm（与图题注一致）
\captionsetup[table]{skip=14pt, margin=3cm}
% 图题注左右各缩进 3cm
\captionsetup[figure]{margin=3cm}
\usepackage{tikz}
\usetikzlibrary{positioning,arrows.meta,shapes.geometric,calc}
% 西文用 Times New Roman（系统字体，XeLaTeX）
\usepackage{fontspec}
\setmainfont{Times New Roman}
% 等宽字体用于代码标识符（fontspec 下 \ttdefault 覆盖与 TU 不兼容）
\setmonofont{Courier New}
% 页边距：四周 2.04cm（1in 减 0.5cm），左右再减 0.5cm -> 1.54cm；上下保持 2.04cm
\usepackage[left=1.54cm, right=1.54cm, top=2.04cm, bottom=2.04cm]{geometry}
\usepackage{setspace}
\onehalfspacing
% 数字编号引用：[1]、[2]……
\usepackage[numbers,square,comma]{natbib}
\setcitestyle{square,comma,numbers}
\usepackage{fancyhdr}
\pagestyle{fancy}
\fancyhf{}
\fancyfoot[C]{\thepage}
\renewcommand{\headrulewidth}{0pt}
\usepackage{hyperref}
\hypersetup{colorlinks=true, linkcolor=black, urlcolor=blue, citecolor=blue,
    pageanchor=false}
% 允许 LaTeX 略微拉伸 glue，避免出现边缘 overfull hbox。
\setlength{\emergencystretch}{2em}
\def\BibTeX{{\rm B\kern-.05em{\sc i\kern-.025em b}\kern-.08em
    T\kern-.1667em\lower.7ex\hbox{E}\kern-.125emX}}

\begin{document}

% ================= 封面 =================
\begin{titlepage}
\thispagestyle{empty}
\centering
% --- 顶部：学校标识 ---
{\large \textbf{香港中文大学}\par}
\vspace{0.25cm}
{\normalsize 信息工程学系，工程学院\par}
\vspace{0.8cm}
{\rule{\textwidth}{0.6pt}\par}
\vfill

% --- 标题区 ---
\begin{minipage}{0.92\textwidth}
\centering
{\normalsize\textbf{暑期科研实习项目}\par}
\vspace{0.9cm}
{\LARGE \textbf{面向 GUI 结构错误修正的\\[0.15cm]异构二分图神经网络}\par}
\vspace{0.6cm}
{\large \textit{面向轻量级视觉语言模型的后修正框架}\par}
\end{minipage}
\vfill

% --- 作者 / 导师并排 ---
\begin{tabular}{@{}l@{\hspace{4.5cm}}l@{}}
\textbf{作者：} & \textbf{导师：} \\
Alex Licheng Xie & Prof.\ Wing Cheong Lau \\
1155246441 & \\
\end{tabular}
\vspace{1.4cm}

% --- 日期 ---
{\normalsize 2026 年 8 月\par}
\vfill
\end{titlepage}

\tableofcontents
\newpage

\begin{abstract}
能够“看”手机屏幕的计算机程序——视障用户辅助工具、自动化 UI 测试器和代替用户操作手机的软件代理——都需要知道屏幕上有什么、在什么位置。轻量级视觉语言模型（VLM）速度快、内存占用低，很适合在手机上直接运行这类任务，但它们会犯错：会漏掉 10--30\% 的可见界面元素，而且对检测到的元素所画的边界框往往与真实位置存在偏移。本报告提出一个后修正框架，在不重新训练 VLM、也不额外添加检测器的情况下修复这些错误。其基本假设是：设计良好的界面遵循空间规则——元素对齐、位于容器之内、间距一致。我们把规则编码为一张异构二分图：一侧是检测到的元素节点，另一侧是元素之间的空间约束节点，边只存在于两类节点之间。GraphSAGE 网络沿这张图传递消息，使每个元素的预测都能参考其空间邻居的信息，然后预测每个元素的坐标修正量、标记被违反的约束、给每条检测打分（真实还是幻觉），并提出 VLM 完全漏掉的元素。在 RICO 数据集（真实安卓截图）上，模型检测违反约束的准确率达到 90--94\%；当大量元素缺失时，补全布局的 IoU 比最近邻基线最高提升 56\%。以 Qwen3-VL Flash 为前端在 200 张真实截图上端到端测试，修正后的流水线找回 106 个此前被漏掉的元素，F1 提升 2.0 个百分点、召回率提升 2.2 个百分点。修正模块运行在 VLM 之后，不修改底层模型，每张截图只增加不到一毫秒的推理时间。
\end{abstract}

\section{引言}
要理解一个智能手机界面，程序首先必须回答一个基本问题：\emph{屏幕上有什么，在哪里？}所有应用——视障用户的屏幕阅读器、自动化 UI 测试器、代替用户操作手机的软件代理——都依赖对界面元素（按钮、文本框、图标、图片）及其位置的准确认知。

\subsection{为什么选择轻量级视觉语言模型}
视觉语言模型（VLM）是能用自然语言描述图像的神经网络；给定一张截图，它们能列出界面元素及其边界框。70 亿参数以上的大型 VLM 在这项任务上相当准确，但对手机而言太慢、太耗内存。30 亿参数以下的轻量级 VLM（如 Qwen3-VL Flash \cite{bai2023qwen} 和 MiniMax-VL-01）足够快，适合端上部署，因此对实时应用很有吸引力。

\subsection{两类系统性失效模式}
在真实截图上的评测表明，轻量级 VLM 会犯两类错误：

\begin{itemize}
\item \textbf{元素遗漏。}10--30\% 的可见元素从未被报告。小图标、分隔线和嵌套容器最常被漏掉。
\item \textbf{位置偏移。}检测到元素周围的边界框可能偏离真实位置数十像素，破坏下游对布局的推理。
\end{itemize}

这些不是随机故障，而是把大模型压缩成小模型的系统性后果。最自然的应对——训练更大的模型——恰恰是端上部署的限制所禁止的。

\subsection{方法概述}
早期实验塑造了本设计。我们最初尝试直接修正坐标，用 MLP 在逐元素特征上训练。结果始终不佳：在模拟高斯噪声上，模型从未超过 no-op 基线；在真实 VLM 输出（Qwen3-VL Plus 和 Flash）上，位置误差本来已经很小，几乎没有可修正的东西。本地 LLaVA-7B 模型每张截图产生的检测太少，图承载不了有用的结构。这些失败指向一个共同原因：单元素特征携带的上下文太少，不足以支撑坐标预测，无论回归器怎么训练。

本报告探索的替代方案是利用结构。GUI 截图按空间规则排布，如对齐、包含、间距和网格归属。这些规则是与像素无关的独立证据，只看像素的 VLM 用不上：如果一行里的元素除了一个以外全部左对齐，那么无论像素长什么样，那个异类很可能位置错了。我们把这些规则形式化为一张图：一侧是元素节点，另一侧是约束节点，边连接元素与其参与的所有约束。图神经网络（GraphSAGE）在这张结构上进行推理，输出修正后的坐标、被违反的约束、逐元素置信度，以及（针对遗漏问题）全新的元素提案。因为修正发生在 VLM 之后，它不修改底层模型：同一个修正网络可以服务于任何输出带边界框元素的 VLM。

本报告其余部分组织如下。第~2 节总结现有方案及其局限。第~3 节详细描述所提方法。第~4 节给出实验与结果。第~5 节讨论局限与未来工作，第~6 节总结。

\section{现有方案及其局限}
\subsection{微调更大的模型}
修复 VLM 错误最直接的办法是在 GUI 数据上微调更大的 VLM（7B+），或在领域特定截图上微调轻量模型。这能提升准确率，但有两个代价：需要带标注的 GUI 数据，而且（对端上场景更根本的是）改进后的模型更大、更慢，或两者兼有。微调还固化了特定的 VLM；更换前端模型意味着全部重新训练。

\subsection{目标检测器级联}
另一种思路是在 VLM 之前或旁边加一个专门的目标检测器（例如 DETR \cite{carion2020detr}）。检测器准确，但要在端上多跑一个网络，增加延迟和内存，还需要自己的训练数据。级联方案把 GUI 当作通用目标检测问题，丢弃了所有设计良好的界面共有的强空间结构。

\subsection{生成式布局模型}
LayoutGAN \cite{li2019layoutgan} 和 LayoutTransformer \cite{gupta2021layouttransformer} 学习布局先验，能\emph{生成}新布局。它们是生成式的：给定随机噪声或部分结构，它们产出新布局。我们的问题不同：我们拿到的是带噪声的布局，要\emph{修复}它，但“学到的布局先验”这一思想是共通的。我们把先验做得显式、可解释（十种空间约束类型），而不是隐含在生成模型里。

\subsection{为什么需要修正网络}
本框架与上述方案有一个根本区别：它从不从零生成像素或检测，也从不修改 VLM。它用结构推理\emph{后处理} VLM 的噪声输出。这一选择源于端上约束：VLM 才是昂贵的组件，任何迫使换更大模型或加检测器的修复都会抵消使用轻量级 VLM 的意义。后处理保持 VLM 不变，只增加一个小的图网络（57K 参数），并且适用于任何输出带边界框元素列表的 VLM。代价是修正受限于 VLM 最初检测到的东西；完全漏掉的元素只能靠补全头找回，这一点在第~4 节评估。

\section{方法}
\subsection{总体流程}
流水线如图~\ref{fig:pipeline} 所示。截图送入轻量级 VLM，得到一份带噪声的检测元素 JSON 列表。我们把该列表转换成异构二分图，跑两跳 GraphSAGE 消息传递，从四个预测头读出修正结果。输出是一份修正后的 JSON：精修过的框、针对漏检元素的新提案、以及逐元素可靠性分数。

\begin{figure}[htbp]
\centering
\resizebox{\textwidth}{!}{%
\begin{tikzpicture}[
    box/.style={draw, rounded corners=2pt, align=center, minimum height=0.9cm,
                font=\small},
    arr/.style={-{Stealth[length=2.2mm]}, thick},
    scale=1.0, transform shape]
  \node[box] (shot) at (0,0) {截图};
  \node[box] (vlm) at (3.0,0) {轻量级\\VLM};
  \node[box, fill=red!8] (noisy) at (6.2,0) {噪声 JSON};
  \node[box, fill=blue!8] (graph) at (9.6,0) {异构\\二分图};
  \node[box, fill=blue!8] (gnn) at (12.9,0) {GraphSAGE};
  \node[box, fill=green!10] (out) at (16.0,0) {修正后 JSON};
  \draw[arr] (shot) -- (vlm);
  \draw[arr] (vlm) -- (noisy);
  \draw[arr] (noisy) -- (graph);
  \draw[arr] (graph) -- (gnn);
  \draw[arr] (gnn) -- (out);
  \node[font=\scriptsize, below=0.28cm of vlm] {检测};
  \node[font=\scriptsize, align=center, below=0.28cm of noisy] {10--30\% 遗漏，\\框偏移};
  \node[font=\scriptsize, below=0.28cm of graph] {约束提取};
  \node[font=\scriptsize, below=0.28cm of gnn] {$\Delta\mathbf{x}$，提案};
\end{tikzpicture}}
\caption{流水线总览：截图、轻量级 VLM、噪声元素列表、二分图、GraphSAGE、修正后 JSON。}
\label{fig:pipeline}
\end{figure}

\subsection{把截图形式化为图}
\subsubsection{元素}
\sloppy
设一张截图产生 $N$ 个检测元素
$\mathcal{E}=\{e_i\}_{i=1}^{N}$，每个元素有一个归一化边界框
$\mathbf{b}_i=(x_1^{(i)},y_1^{(i)},x_2^{(i)},y_2^{(i)})\in[0,1]^4$（框坐标，缩放至单位正方形）和一个来自元素分类体系的类型标签
$t_i\in\mathcal{T}$（按钮、文本框、图标、图片等）。

\subsubsection{空间约束}
由这些元素提取 $M$ 条空间约束
$\mathcal{C}=\{c_j\}_{j=1}^{M}$，每条约束是元素之间的一种类型化关系。表~\ref{tab:constraints} 列出十种约束类型。每条约束 $c_j$ 有类型 $\tau_j$、源索引集合
$S_j\subseteq\{1,\dots,N\}$ 和目标索引集合
$T_j\subseteq\{1,\dots,N\}$；当谓词
$P_\tau(\mathbf{b}_{S_j},\mathbf{b}_{T_j})$ 在容差
$\varepsilon$ 内成立时，该约束存在。

\begin{table}[htbp]
\caption{从布局中提取的十种空间约束类型。}
\label{tab:constraints}
\centering
\begin{tabularx}{\textwidth}{lXX}
\toprule
\textbf{类型} & \textbf{谓词} & \textbf{示例} \\
\midrule
\texttt{ALIGN\_LEFT} & $|x_1-x_1'|<\varepsilon$ & 按钮共享左边缘 \\
\texttt{ALIGN\_RIGHT} & $|x_2-x_2'|<\varepsilon$ & 右边缘对齐 \\
\texttt{ALIGN\_TOP} & $|y_1-y_1'|<\varepsilon$ & 上边缘对齐 \\
\texttt{ALIGN\_BOTTOM} & $|y_2-y_2'|<\varepsilon$ & 下边缘对齐 \\
\texttt{CENTER\_X} & 中心距 $<\varepsilon$ & 水平居中 \\
\texttt{CENTER\_Y} & 中心距 $<\varepsilon$ & 垂直居中 \\
\texttt{SPACING} & 间距一致 & 列表项等距排布 \\
\texttt{CONTAINMENT} & 一个框在另一个框内 & 图标位于容器内 \\
\texttt{GRID} & 行/列归属 & 网格布局检测 \\
\texttt{SAME\_SIZE} & 宽/高相近 & 卡片尺寸统一 \\
\bottomrule
\end{tabularx}
\end{table}

\subsubsection{二分图}
GUI 结构是一张\emph{异构二分图}
\begin{equation}
G=(\mathcal{V},\mathcal{E}_{\text{edge}},\phi,\psi),
\label{eq:graph}
\end{equation}
其中 $\mathcal{V}=\mathcal{E}\cup\mathcal{C}$ 是节点集，分成元素节点 $V_e$（$|V_e|=N$）和约束节点 $V_c$（$|V_c|=M$）；$\phi:\mathcal{V}\to\{0,1\}$ 标记每个节点所属的分区；边只存在于两个分区之间；当且仅当 $e_i\in S_j\cup T_j$ 时存在边 $(e_i,c_j)$ 连接元素 $e_i$ 与约束 $c_j$；$\psi$ 用元素与约束子空间之间的归一化距离为每条边加权。不存在元素--元素或约束--约束边。图~\ref{fig:arch} 展示了这一结构。

\begin{figure}[htbp]
\centering
\begin{tikzpicture}[
    elem/.style={draw, rounded corners=2pt, fill=blue!10, minimum width=1.9cm,
                 minimum height=0.55cm, font=\footnotesize},
    cons/.style={draw, rounded corners=2pt, fill=orange!15, minimum width=1.9cm,
                 minimum height=0.55cm, font=\footnotesize},
    edg/.style={gray!60, thick},
    scale=1.0, transform shape]
  % 元素节点（左）
  \node[elem] (e1) at (0,2.6) {$e_1$: 按钮};
  \node[elem] (e2) at (0,1.6) {$e_2$: 文本};
  \node[elem] (e3) at (0,0.6) {$e_3$: 图标};
  \node[elem] (e4) at (0,-0.4) {$e_4$: 输入框};
  \node[elem] (e5) at (0,-1.4) {$e_5$: 图片};
  % 约束节点（右）
  \node[cons] (c1) at (5.2,2.0) {$c_1$: \texttt{ALIGN\_LEFT}};
  \node[cons] (c2) at (5.2,0.9) {$c_2$: \texttt{CONTAINMENT}};
  \node[cons] (c3) at (5.2,-0.2) {$c_3$: \texttt{SPACING}};
  \node[cons] (c4) at (5.2,-1.3) {$c_4$: \texttt{SAME\_SIZE}};
  % 边：只跨分区
  \draw[edg] (e1) -- (c1);
  \draw[edg] (e1) -- (c2);
  \draw[edg] (e2) -- (c1);
  \draw[edg] (e2) -- (c3);
  \draw[edg] (e3) -- (c2);
  \draw[edg] (e3) -- (c4);
  \draw[edg] (e4) -- (c3);
  \draw[edg] (e5) -- (c4);
  \draw[edg] (e5) -- (c2);
  % 分区标签
  \node[font=\small\itshape, above=0.25cm of e1] {元素节点 $V_e$};
  \node[font=\small\itshape, above=0.25cm of c1] {约束节点 $V_c$};
\end{tikzpicture}
\caption{异构二分图：元素节点（左）只连接约束节点（右）。任意两个元素从不直接相连；它们通过共享约束通信。}
\label{fig:arch}
\end{figure}

\subsection{消息传递}
图神经网络沿边传播信息。我们使用两跳交替的 GraphSAGE 卷积 \cite{hamilton2017sage}，用均值操作聚合邻居特征。

\subsubsection{第一跳：元素到约束}
每个约束节点收集它所涉及的元素的聚合信息：
\begin{equation}
\mathbf{h}_{c_j}^{(1)}=\sigma\!\left(\mathbf{W}_1\cdot
\operatorname{MEAN}\!\left(\left\{\mathbf{h}_{e_i}^{(0)}:
(e_i,c_j)\in\mathcal{E}_{\text{edge}}\right\}\right)+\mathbf{b}_1\right),
\label{eq:hop1}
\end{equation}
其中 $\mathbf{h}_{e_i}^{(0)}$ 是初始元素特征向量，$\sigma$ 是 ReLU 激活。直观地说，约束节点现在“知道”共享它的元素的总体状态。

\subsubsection{第二跳：约束到元素}
每个元素节点把更新后的约束状态聚合回来：
\begin{equation}
\mathbf{h}_{e_i}^{(2)}=\sigma\!\left(\mathbf{W}_2\cdot
\operatorname{MEAN}\!\left(\left\{\mathbf{h}_{c_j}^{(1)}:
(e_i,c_j)\in\mathcal{E}_{\text{edge}}\right\}\right)+\mathbf{b}_2\right).
\label{eq:hop2}
\end{equation}
这一跳之后，每个元素的表示都携带了与它共享约束的所有元素的信息，即它的“空间邻域”。不共享约束的元素之间从不交换消息，这保持了强归纳偏置，也让计算保持局部性。

\subsection{节点特征}
每个元素的初始特征 $\mathbf{h}_{e_i}^{(0)}$ 是 5 维的：四个归一化框坐标加上归一化面积 $a_i=(x_2-x_1)(y_2-y_1)$。可选地拼接冻结的视觉特征 $\mathbf{v}_i\in\mathbb{R}^{d_v}$（192 维 ViT-Tiny 或 768 维 DINOv2 \cite{oquab2024dinov2}）。约束节点特征把约束类型编码为 10 维 one-hot 向量，加上参与元素的空间统计量（平均两两距离、包含重叠比、对齐残差）。

\subsection{四个预测头}
消息传递之后，四个头读出结果。总训练目标是加权和
\begin{equation}
\mathcal{L}=w_c\,\mathcal{L}_{\text{coord}}+w_v\,\mathcal{L}_{\text{vio}}
+w_e\,\mathcal{L}_{\text{exist}}+w_p\,\mathcal{L}_{\text{prop}},
\label{eq:loss}
\end{equation}
默认权重 $w_c{=}1.0$，$w_v{=}0.5$，$w_e{=}0.5$，$w_p{=}0.5$。

\subsubsection{坐标修正}
第一个头预测每个元素的修正向量
$\Delta\mathbf{x}_i=\operatorname{MLP}_{\text{coord}}(\mathbf{h}_{e_i}^{(2)})$，
即 VLM 的框应移动的量，用 smooth L1 损失优化：
\begin{equation}
\mathcal{L}_{\text{coord}}=\frac{1}{N}\sum_{i=1}^{N}
\operatorname{smooth}_{L_1}(\Delta\mathbf{x}_i-\Delta\mathbf{x}_i^*),
\label{eq:lcoord}
\end{equation}
其中 $\Delta\mathbf{x}_i^*$ 是真实修正量。

\subsubsection{违规检测}
第二个头对每条约束是否被违反（被破坏）分类，$\hat{v}_j=\sigma(\operatorname{MLP}_{\text{vio}}(\mathbf{h}_{c_j}^{(1)}))$：
\begin{equation}
\mathcal{L}_{\text{vio}}=-\frac{1}{M}\sum_{j=1}^{M}\bigl[
v_j^*\log\hat{v}_j+(1-v_j^*)\log(1-\hat{v}_j)\bigr],
\label{eq:lvio}
\end{equation}
其中 $v_j^*\in\{0,1\}$ 表示约束 $c_j$ 在真实布局中是否确实被违反。正是这个头让模型“理解”布局规则：它必须仅凭图学会破坏的对齐或包含长什么样。

\subsubsection{存在性评分}
第三个头给每个检测元素打分，判断它是真实的还是幻觉，在元素嵌入上用二分类器预测 $\hat{e}_i$，用二元交叉熵训练：
\begin{equation}
\mathcal{L}_{\text{exist}}=-\frac{1}{N}\sum_{i=1}^{N}\bigl[
e_i^*\log\hat{e}_i+(1-e_i^*)\log(1-\hat{e}_i)\bigr],
\label{eq:lexist}
\end{equation}
其中 $e_i^*$ 是真实存在性标签。下游可以用这个分数过滤假阳性，或按可靠性给元素排序。

\subsubsection{元素补全}
第四个头解决遗漏问题，提出\emph{缺失的}元素。训练时，我们随机丢弃 $\rho\in[0.2,0.8]$ 比例的真实元素。引用被丢弃元素的约束现在有了一个“洞”（一条悬空边），这正是缺失元素在图中留下的签名。该头从聚合的约束嵌入预测缺失元素的框和类型：
\begin{equation}
\hat{\mathbf{b}}_k,\hat{t}_k=
\operatorname{MLP}_{\text{proposal}}(\mathbf{h}_{c_j}^{(1)}),
\qquad
\mathcal{L}_{\text{prop}}=\frac{1}{K}\sum_{k=1}^{K}\bigl[
\mathcal{L}_{\text{IoU}}(\hat{\mathbf{b}}_k,\mathbf{b}_k^*)
+\alpha\,\mathrm{CE}(\hat{t}_k,t_k^*)\bigr],
\label{eq:lprop}
\end{equation}
其中 $K$ 是目标缺失的违规约束数量，$\mathcal{L}_{\text{IoU}}$ 是基于 IoU 的框损失。注意训练信号完全来自对真实布局的掩蔽；补全头是\emph{自监督}的，不需要额外的人工标注。

\section{实验与结果}
\subsection{实验设置}
所有模型使用 AdamW 优化器和余弦退火调度，隐藏维度 128，两层二分 GraphSAGE；报告的结果用 5 个随机种子训练。我们在 RICO 数据集 \cite{dek.a17rico}（真实安卓截图）和 ScreenSpot \cite{cheng2024seclick} 上评估。实现使用 PyTorch \cite{paszke2019pytorch}、PyTorch Geometric \cite{fey2019pyg}、NumPy \cite{harris2020numpy}、SciPy \cite{virtanen2020scipy} 和 Matplotlib \cite{hunter2007matplotlib}；前端 VLM 是 Qwen3-VL Flash \cite{bai2023qwen}。

\subsection{约束类型消融}
哪些空间规则最重要？我们在违规检测任务上消融十种约束类型（表~\ref{tab:ablation}、图~\ref{fig:ablation}）。移除包含约束造成最大准确率下降（$-1.9$pp），证实父子包含是最强的结构信号。移除\emph{全部}对齐类型反而把准确率\emph{提高到} 93.9\%：每张图上的约束少得多，剩下的违规更容易分类。只保留对齐类型损失最大（$-3.0$pp）：仅靠对齐支撑不起违规检测。

\begin{table}[htbp]
\caption{违规检测上的约束类型消融（n=500，drop=0.6）。}
\label{tab:ablation}
\centering
\begin{tabularx}{\textwidth}{Xc}
\toprule
\textbf{约束集合} & \textbf{违规准确率} \\
\midrule
全部 10 类（对照组） & \textbf{0.908} \\
移除 \texttt{CONTAINMENT} & 0.889（$-1.9$pp） \\
移除 \texttt{SPACING} & 0.903（$-0.5$pp） \\
移除 \texttt{GRID} & 0.916（$+0.8$pp） \\
移除全部 \texttt{ALIGNMENT} & 0.939（$+3.1$pp） \\
仅保留 \texttt{ALIGNMENT} & 0.878（$-3.0$pp） \\
\bottomrule
\end{tabularx}
\end{table}

\begin{figure}[htbp]
\centering
\begin{tikzpicture}[scale=1.0, transform shape]
  % 水平条形图：长度 = (值 - 0.84) * 60 cm
  \def\sc{60}
  \def\base{0.0}
  % 6 行的 y 坐标（自上而下），条高 0.35；按值降序排列
  % 移除全部 ALIGNMENT 0.939 -> 5.94cm
  \fill[gray!50] (\base,5.0) rectangle ({(\base+(0.939-0.84)*\sc)},5.35);
  \node[anchor=east, font=\footnotesize] at (\base-0.1,5.175) {移除全部 \texttt{ALIGNMENT}};
  \node[font=\footnotesize, anchor=west] at ({(\base+(0.939-0.84)*\sc+0.1)},5.175) {0.939};
  % 移除 GRID 0.916 -> 4.56cm
  \fill[gray!50] (\base,4.0) rectangle ({(\base+(0.916-0.84)*\sc)},4.35);
  \node[anchor=east, font=\footnotesize] at (\base-0.1,4.175) {移除 \texttt{GRID}};
  \node[font=\footnotesize, anchor=west] at ({(\base+(0.916-0.84)*\sc+0.1)},4.175) {0.916};
  % 全部 10 类（对照组）0.908 -> 4.08cm
  \fill[blue!60] (\base,3.0) rectangle ({(\base+(0.908-0.84)*\sc)},3.35);
  \node[anchor=east, font=\footnotesize] at (\base-0.1,3.175) {全部 10 类};
  \node[font=\footnotesize, anchor=west] at ({(\base+(0.908-0.84)*\sc+0.1)},3.175) {0.908};
  % 移除 SPACING 0.903 -> 3.78cm
  \fill[gray!50] (\base,2.0) rectangle ({(\base+(0.903-0.84)*\sc)},2.35);
  \node[anchor=east, font=\footnotesize] at (\base-0.1,2.175) {移除 \texttt{SPACING}};
  \node[font=\footnotesize, anchor=west] at ({(\base+(0.903-0.84)*\sc+0.1)},2.175) {0.903};
  % 移除 CONTAINMENT 0.889 -> 2.94cm
  \fill[gray!50] (\base,1.0) rectangle ({(\base+(0.889-0.84)*\sc)},1.35);
  \node[anchor=east, font=\footnotesize] at (\base-0.1,1.175) {移除 \texttt{CONTAINMENT}};
  \node[font=\footnotesize, anchor=west] at ({(\base+(0.889-0.84)*\sc+0.1)},1.175) {0.889};
  % 仅保留 ALIGNMENT 0.878 -> 2.28cm
  \fill[gray!50] (\base,0.0) rectangle ({(\base+(0.878-0.84)*\sc)},0.35);
  \node[anchor=east, font=\footnotesize] at (\base-0.1,0.175) {仅保留 \texttt{ALIGNMENT}};
  \node[font=\footnotesize, anchor=west] at ({(\base+(0.878-0.84)*\sc+0.1)},0.175) {0.878};
  % 坐标轴
  \draw[thick,-{Stealth}] (\base,5.75) -- (\base,-0.5);
  \draw[thick,-{Stealth}] (\base,-0.65) -- (6.6,-0.65);
  \node[font=\footnotesize] at (3.3,-1.5) {违规准确率};
  \foreach \v in {0.86,0.88,0.90,0.92,0.94} {
    \draw ({(\v-0.84)*\sc},-0.65) -- ++(0,-0.08) node[below, font=\tiny] {\v};
  }
\end{tikzpicture}
\caption{各约束集合下的违规准确率。\texttt{CONTAINMENT} 是最关键的单类约束。}
\label{fig:ablation}
\end{figure}

\subsection{训练目标消融}
我们比较联合训练与单目标训练（5 个种子，图~\ref{fig:phase9}）。联合训练取得最佳平衡：违规准确率 0.876、提案 MSE 0.051。仅违规训练达到更高的违规准确率（0.898），但提案 MSE 很差（0.116）；仅提案训练提案 MSE 很好（0.051），但违规准确率接近随机（0.489）。联合目标用一点违规准确率换取了很大的提案质量提升，这对补全很重要。

\begin{figure}[htbp]
\centering
\begin{tikzpicture}[scale=1.0, transform shape]
  % 左面板：违规准确率（联合 0.876、仅违规 0.898、仅提案 0.489）
  % 比例：1.0 = 4cm，y = 值*4
  \def\sc{4}
  \begin{scope}[shift={(0,0)}]
    \node[font=\footnotesize] at (2.4,4.5) {违规头};
    % 联合
    \fill[blue!60] (0.6,0) rectangle (1.6,{0.876*\sc});
    % 仅违规
    \fill[blue!60] (2.0,0) rectangle (3.0,{0.898*\sc});
    % 仅提案
    \fill[blue!60] (3.4,0) rectangle (4.4,{0.489*\sc});
    % 误差条带帽（std: 0.019, 0.018, 0.071）；标签在帽上方
    \draw (1.1,{0.876*\sc+0.05}) -- (1.1,{0.876*\sc+0.05+0.019*\sc});
    \draw (1.02,{0.876*\sc+0.05+0.019*\sc}) -- (1.18,{0.876*\sc+0.05+0.019*\sc});
    \node[font=\scriptsize] at (1.1,{0.876*\sc+0.05+0.019*\sc+0.18}) {0.876};
    \draw (2.5,{0.898*\sc+0.05}) -- (2.5,{0.898*\sc+0.05+0.018*\sc});
    \draw (2.42,{0.898*\sc+0.05+0.018*\sc}) -- (2.58,{0.898*\sc+0.05+0.018*\sc});
    \node[font=\scriptsize] at (2.5,{0.898*\sc+0.05+0.018*\sc+0.18}) {0.898};
    \draw (3.9,{0.489*\sc+0.05}) -- (3.9,{0.489*\sc+0.05+0.071*\sc});
    \draw (3.82,{0.489*\sc+0.05+0.071*\sc}) -- (3.98,{0.489*\sc+0.05+0.071*\sc});
    \node[font=\scriptsize] at (3.9,{0.489*\sc+0.05+0.071*\sc+0.18}) {0.489};
    \draw[thick,-{Stealth}] (0,-0.15) -- (0,4.1);
    \draw[thick,-{Stealth}] (-0.15,0) -- (4.6,0);
    \node[font=\scriptsize, rotate=45, anchor=east] at (1.1,-0.45) {联合};
    \node[font=\scriptsize, rotate=45, anchor=east] at (2.5,-0.45) {仅违规};
    \node[font=\scriptsize, rotate=45, anchor=east] at (3.9,-0.45) {仅提案};
  \end{scope}
  % 右面板：提案 MSE（联合 0.051、仅违规 0.116、仅提案 0.051）
  % MSE 比例：1.0 = 30cm，y = 值*30（0.116 -> 3.48cm，与左面板一致）
  \def\scm{30}
  \begin{scope}[shift={(6.6,0)}]
    \node[font=\footnotesize] at (2.4,4.5) {提案头};
    \fill[blue!60] (0.6,0) rectangle (1.6,{0.051*\scm});
    \fill[blue!60] (2.0,0) rectangle (3.0,{0.116*\scm});
    \fill[blue!60] (3.4,0) rectangle (4.4,{0.051*\scm});
    % 误差条带帽（std: 0.009, 0.009, 0.010）*30；标签在帽上方
    \draw (1.1,{0.051*\scm+0.05}) -- (1.1,{0.051*\scm+0.05+0.009*\scm});
    \draw (1.02,{0.051*\scm+0.05+0.009*\scm}) -- (1.18,{0.051*\scm+0.05+0.009*\scm});
    \node[font=\scriptsize] at (1.1,{0.051*\scm+0.05+0.009*\scm+0.18}) {0.051};
    \draw (2.5,{0.116*\scm+0.05}) -- (2.5,{0.116*\scm+0.05+0.009*\scm});
    \draw (2.42,{0.116*\scm+0.05+0.009*\scm}) -- (2.58,{0.116*\scm+0.05+0.009*\scm});
    \node[font=\scriptsize] at (2.5,{0.116*\scm+0.05+0.009*\scm+0.18}) {0.116};
    \draw (3.9,{0.051*\scm+0.05}) -- (3.9,{0.051*\scm+0.05+0.010*\scm});
    \draw (3.82,{0.051*\scm+0.05+0.010*\scm}) -- (3.98,{0.051*\scm+0.05+0.010*\scm});
    \node[font=\scriptsize] at (3.9,{0.051*\scm+0.05+0.010*\scm+0.18}) {0.051};
    \draw[thick,-{Stealth}] (0,-0.15) -- (0,4.1);
    \draw[thick,-{Stealth}] (-0.15,0) -- (4.6,0);
    \node[font=\scriptsize, rotate=45, anchor=east] at (1.1,-0.45) {联合};
    \node[font=\scriptsize, rotate=45, anchor=east] at (2.5,-0.45) {仅违规};
    \node[font=\scriptsize, rotate=45, anchor=east] at (3.9,-0.45) {仅提案};
  \end{scope}
\end{tikzpicture}
\caption{训练目标消融（5 个种子，均值 $\pm$ 标准差）：联合训练兼顾两个头，单目标训练则损害另一个头。}
\label{fig:phase9}
\end{figure}

\subsection{元素补全}
我们把 GNN 提案头与最近邻基线比较，后者复制最近存活元素的框（图~\ref{fig:completion}）。在低缺失比例（0.2--0.4）下基线占优，因为附近有很多存活元素。在高缺失比例（0.6--0.8）下 GNN 占优：drop 0.6 时 GNN 的 IoU 为 0.122，NN 为 0.088（$+39\%$）；drop 0.8 时为 0.097 对 0.062（$+56\%$）。存活元素越少，结构推理比邻近性越重要，这与模型利用结构先验而非简单插值是一致的。

\begin{figure}[htbp]
\centering
\begin{tikzpicture}[scale=1.0, transform shape]
  % 分组柱状图：4 个缺失比例 x {GNN, NN}；y = 值 * 22 cm（最大约 0.204 -> 4.5cm）
  \def\sc{22}
  % drop 0.2: GNN 0.047+/-0.009, NN 0.057+/-0.031
  % drop 0.4: GNN 0.079+/-0.014, NN 0.110+/-0.013
  % drop 0.6: GNN 0.122+/-0.082, NN 0.088+/-0.011
  % drop 0.8: GNN 0.097+/-0.047, NN 0.062+/-0.012
  % x 坐标：第 i 组中心在 1+2i，柱子在中线左右各 0.3
  \begin{scope}
    % 图例（顶部中央，轴上方）
    \fill[blue!60] (2.4,5.3) rectangle (2.8,5.55);
    \node[anchor=west, font=\scriptsize] at (2.85,5.42) {GNN（本方法）};
    \fill[gray!60] (4.9,5.3) rectangle (5.3,5.55);
    \node[anchor=west, font=\scriptsize] at (5.35,5.42) {NN 基线};
    % 第 1 组：drop 0.2
    \fill[blue!60] (0.7,0) rectangle (1.3,{0.047*\sc});
    \draw (1.0,{0.047*\sc+0.04}) -- (1.0,{0.047*\sc+0.04+0.009*\sc});
    \draw (0.92,{0.047*\sc+0.04+0.009*\sc}) -- (1.08,{0.047*\sc+0.04+0.009*\sc});
    \fill[gray!60] (1.7,0) rectangle (2.3,{0.057*\sc});
    \draw (2.0,{0.057*\sc+0.04}) -- (2.0,{0.057*\sc+0.04+0.031*\sc});
    \draw (1.92,{0.057*\sc+0.04+0.031*\sc}) -- (2.08,{0.057*\sc+0.04+0.031*\sc});
    \node[font=\scriptsize] at (1.5,-0.28) {0.2};
    % 第 2 组：drop 0.4
    \fill[blue!60] (2.7,0) rectangle (3.3,{0.079*\sc});
    \draw (3.0,{0.079*\sc+0.04}) -- (3.0,{0.079*\sc+0.04+0.014*\sc});
    \draw (2.92,{0.079*\sc+0.04+0.014*\sc}) -- (3.08,{0.079*\sc+0.04+0.014*\sc});
    \fill[gray!60] (3.7,0) rectangle (4.3,{0.110*\sc});
    \draw (4.0,{0.110*\sc+0.04}) -- (4.0,{0.110*\sc+0.04+0.013*\sc});
    \draw (3.92,{0.110*\sc+0.04+0.013*\sc}) -- (4.08,{0.110*\sc+0.04+0.013*\sc});
    \node[font=\scriptsize] at (3.5,-0.28) {0.4};
    % 第 3 组：drop 0.6
    \fill[blue!60] (4.7,0) rectangle (5.3,{0.122*\sc});
    \draw (5.0,{0.122*\sc+0.04}) -- (5.0,{0.122*\sc+0.04+0.082*\sc});
    \draw (4.92,{0.122*\sc+0.04+0.082*\sc}) -- (5.08,{0.122*\sc+0.04+0.082*\sc});
    \fill[gray!60] (5.7,0) rectangle (6.3,{0.088*\sc});
    \draw (6.0,{0.088*\sc+0.04}) -- (6.0,{0.088*\sc+0.04+0.011*\sc});
    \draw (5.92,{0.088*\sc+0.04+0.011*\sc}) -- (6.08,{0.088*\sc+0.04+0.011*\sc});
    \node[font=\scriptsize] at (5.5,-0.28) {0.6};
    % 第 4 组：drop 0.8
    \fill[blue!60] (6.7,0) rectangle (7.3,{0.097*\sc});
    \draw (7.0,{0.097*\sc+0.04}) -- (7.0,{0.097*\sc+0.04+0.047*\sc});
    \draw (6.92,{0.097*\sc+0.04+0.047*\sc}) -- (7.08,{0.097*\sc+0.04+0.047*\sc});
    \fill[gray!60] (7.7,0) rectangle (8.3,{0.062*\sc});
    \draw (8.0,{0.062*\sc+0.04}) -- (8.0,{0.062*\sc+0.04+0.012*\sc});
    \draw (7.92,{0.062*\sc+0.04+0.012*\sc}) -- (8.08,{0.062*\sc+0.04+0.012*\sc});
    \node[font=\scriptsize] at (7.5,-0.28) {0.8};
    % 坐标轴
    \draw[thick,-{Stealth}] (0,-0.12) -- (0,4.9);
    \draw[thick,-{Stealth}] (-0.15,0) -- (8.7,0);
    \node[font=\footnotesize] at (4.3,-0.75) {缺失比例};
    \node[font=\footnotesize, rotate=90] at (-1.45,2.4) {提案 IoU};
    \foreach \v in {0.05,0.10,0.15,0.20} {
      \draw (-0.1,{\v*\sc}) -- (0.1,{\v*\sc});
      \node[anchor=east, font=\tiny] at (-0.15,{\v*\sc}) {\v};
    }
  \end{scope}
\end{tikzpicture}
\caption{元素补全 IoU 随缺失比例的变化（均值 $\pm$ 标准差，2 个种子）。当大量结构缺失时（drop 0.6--0.8），GNN 超过最近邻基线。}
\label{fig:completion}
\end{figure}

\subsection{真实 VLM 输出的端到端实验}
我们把训练好的模型部署在 Qwen3-VL Flash 之后，在 200 张真实 RICO 截图（4789 个真实元素、2947 条 VLM 预测）上测试。GNN 接收 VLM 的噪声元素，检测被违反的约束并提出缺失元素；提案用非极大值抑制合并。我们用中心距离匹配、阈值 0.1 对真实标注评估（表~\ref{tab:real}、图~\ref{fig:realvlm}）。

\begin{table}[htbp]
\caption{200 张真实 RICO 截图上的端到端结果（合并统计）。Checkpoint 采用形状过滤加载；数字对应正确加载的权重。}
\label{tab:real}
\centering
\begin{tabularx}{\textwidth}{Xccc}
\toprule
\textbf{指标} & \textbf{仅 VLM} & \textbf{VLM + GNN} & \textbf{$\Delta$} \\
\midrule
精确率 & 0.382 & 0.393 & $+1.1$pp \\
召回率 & 0.235 & 0.257 & $+2.2$pp \\
F1 & 0.291 & 0.311 & $+2.0$pp \\
TP / FP / FN & 1126/1821/3663 & 1232/1905/3557 & $+106$ TP \\
\bottomrule
\end{tabularx}
\end{table}

\begin{figure}[htbp]
\centering
\begin{tikzpicture}[scale=1.0, transform shape]
  % 修正前后：精确率 0.382->0.393，召回率 0.235->0.257，F1 0.291->0.311
  % y = 值 * 14 cm（0.393 -> 5.5cm）；轴上界 6.2 为标签留空间
  \def\sc{14}
  % 图例（顶部中央，轴上方）
  \fill[gray!60] (2.1,6.0) rectangle (2.5,6.25);
  \node[anchor=west, font=\scriptsize] at (2.55,6.12) {仅 VLM};
  \fill[blue!60] (4.3,6.0) rectangle (4.7,6.25);
  \node[anchor=west, font=\scriptsize] at (4.75,6.12) {VLM + GNN};
  % 精确率
  \fill[gray!60] (0.7,0) rectangle (1.3,{0.382*\sc});
  \node[font=\scriptsize] at (1.0,{0.382*\sc+0.12}) {0.382};
  \fill[blue!60] (1.7,0) rectangle (2.3,{0.393*\sc});
  \node[font=\scriptsize] at (2.0,{0.393*\sc+0.12}) {0.393};
  \node[font=\scriptsize] at (1.5,-0.28) {精确率};
  % 召回率
  \fill[gray!60] (2.7,0) rectangle (3.3,{0.235*\sc});
  \node[font=\scriptsize] at (3.0,{0.235*\sc+0.12}) {0.235};
  \fill[blue!60] (3.7,0) rectangle (4.3,{0.257*\sc});
  \node[font=\scriptsize] at (4.0,{0.257*\sc+0.12}) {0.257};
  \node[font=\scriptsize] at (3.5,-0.28) {召回率};
  % F1
  \fill[gray!60] (4.7,0) rectangle (5.3,{0.291*\sc});
  \node[font=\scriptsize] at (5.0,{0.291*\sc+0.12}) {0.291};
  \fill[blue!60] (5.7,0) rectangle (6.3,{0.311*\sc});
  \node[font=\scriptsize] at (6.0,{0.311*\sc+0.12}) {0.311};
  \node[font=\scriptsize] at (5.5,-0.28) {F1};
  % 坐标轴
  \draw[thick,-{Stealth}] (0,-0.12) -- (0,5.8);
  \draw[thick,-{Stealth}] (-0.15,0) -- (6.9,0);
  \node[font=\footnotesize, rotate=90] at (-0.95,2.7) {分数};
  \foreach \v in {0.1,0.2,0.3,0.4} {
    \draw (-0.1,{\v*\sc}) -- (0.1,{\v*\sc});
    \node[anchor=east, font=\tiny] at (-0.15,{\v*\sc}) {\v};
  }
\end{tikzpicture}
\caption{端到端对比：仅 VLM 与 VLM + GNN（修正后的 checkpoint 加载）。召回率和 F1 提升，精确率也略有上升。}
\label{fig:realvlm}
\end{figure}

修正后的流水线找回 106 个此前被漏掉的元素（$+106$ TP、$+2.2$pp 召回率），代价是 84 个额外的假阳性，F1 提升 2.0 个百分点。精确率也略有提升（$+1.1$pp），这与早先用错误加载的 checkpoint 报告的结果不同，是一个细微但重要的评测陷阱。早先的一次评测在非严格权重加载下加载了隐藏维度不匹配的 checkpoint，静默丢弃了 89\% 的权重，从近乎随机的权重中产生了一个虚假的 $+2.9$pp F1 提升。这里报告的所有结果都使用形状过滤的 checkpoint 加载——只保留形状匹配的权重——我们建议把它作为标准做法。

\subsection{性能}
修正网络很小（57K 参数）。在一台仅 CPU 的 M3 MacBook Pro 上，每张截图的图构建约 5~ms，推理约 0.53~ms（p50），与每张图约两秒的 VLM 推理相比可以忽略。在 CPU 上训练 2,000 个样本 $\times$ 50 个 epoch 不到五分钟完成。

\section{讨论}
\subsection{结构推理何时有效}
纵观实验，结构推理在 VLM 输出稀疏（大量元素缺失）时帮助最大，在 VLM 已经准确时帮助较小。这呼应了最初促使图形式化的一连串负面结果。在模拟高斯噪声上，模型从未超过 no-op 基线；Qwen3-VL Plus 和 Flash 产生的位置误差在 0.01 左右或更小，没什么可修正；LLaVA-7B 每张截图只检测到几个元素，图承载不了多少结构；重加权的仅存在性目标也没有假阳性可过滤。只有当相当比例的元素缺失时，约束图才包含逐元素回归器无法触及的信息。对真实截图上的轻量级 VLM（漏掉 10--30\% 的元素）来说，这一情形正是实际应用相关的。

\subsection{局限性}
主要局限如下。第一，\emph{类型预测较弱}：仅凭约束上下文预测元素的语义类型，准确率封顶在 62\%；约束携带的是空间信息而非语义信息。第二，\emph{跨域迁移有限}：置信度模型的 AUROC 从 RICO 上的 0.703 降到 ScreenSpot 上的 0.554，说明在手机屏幕上学到的假阳性模式不能完全迁移到混合的移动/PC/网页布局。第三，\emph{密集工具栏很难}：含 8 个以上小图标的工具栏会形成密集的约束图，单个元素难以区分；而自由形式布局（地图、绘图）几乎没有可依据的约束。

\subsection{未来工作}
视觉特征融合（交叉注意力）是一个候选下一步（早期实验显示提案 MSE 改善 18--22\%），对 ScreenSpot 的领域自适应也是。给图加上时间上下文，也许能帮模型区分缺失元素与被 UI 状态隐藏的元素。

\section{结论}
本报告提出一个用于修正轻量级 VLM 噪声 GUI 元素预测的异构二分 GraphSAGE 框架。十种空间约束类型被编码为二分图，元素只通过共享约束交换消息；多任务目标把坐标修正、违规检测、存在性评分和自监督元素补全耦合在一起。在 RICO 上，模型达到 90--94\% 的违规检测准确率，在高缺失比例下补全 IoU 比最近邻基线提升 $+56\%$，并在 200 张真实 Qwen3-VL Flash 截图（形状过滤的 checkpoint 加载）上取得 $+2.0$pp 的 F1 提升。

结果是令人鼓舞的，但也有边界。增益出现在 VLM 输出稀疏时，且框架继承 VLM 最初的检测结果；它不修复纯语义错误，置信度分数也只能部分跨域迁移。这些边界，连同讨论中报告的负面结果，说明结构后修正更适合被看作 GUI 理解流水线的一个组件，而不是完整方案。该方法能否从移动布局推广到桌面和网页界面，仍有待测试。

% 参考文献另起一页
\clearpage
\bibliographystyle{unsrtnat}
\bibliography{references}

\end{document}
